\emph{is of finite type.}

Indeed, by 3.4 (i), one may evidently suppose @@M0@@ reduced and irreducible. Then there exists a nonempty open subset @@M1@@ of @@M2@@ such that @@M3@@ and @@M4@@ are representable and the restriction of @@M5@@ is quasi-affine (XII 1.1), hence of finite type by 3.2 (i). Endowing @@M6@@ with the induced reduced structure, one may suppose that the morphism @@M7@@ induced by @@M8@@ is of finite type by noetherian induction\,; this gives the assertion, by 3.4 (i) applied to @@M9@@.

\sgaStableTarget{sga6.e13.thm.3-6}\textbf{Theorem 3.6.} \emph{Let @@M10@@ be a noetherian scheme, @@M11@@ a proper @@M12@@-scheme and @@M13@@ the Picard functor. Then, for every integer @@M14@@, the @@M15@@-th power homomorphism}
@@M16@@
\emph{is a morphism of finite type.}

Indeed, let @@M17@@ be a Chow presentation: @@M18@@ is projective over @@M19@@ and @@M20@@ is surjective. If the theorem is admitted for @@M21@@, then the theorem for @@M22@@ follows by 3.5 and 3.4 (ii). One may therefore suppose @@M23@@ projective.

There exist a nonempty open subset @@M24@@ of @@M25@@, an integral scheme @@M26@@ and a finite surjective morphism @@M27@@, such that every irreducible component @@M28@@ of @@M29@@, endowed with the induced reduced structure, is integral over @@M30@@ (XII 1.3). Endowing @@M31@@ with the induced reduced structure, one may admit the theorem for @@M32@@ by noetherian induction. If it is admitted for all the @@M33@@, it follows for @@M34@@ by the
